% ------------------ INICIO DE LOS PSEUDOCODIGOS -------------------
\begin{itemize}
    \item \textbf{Pseudocódigos}
\end{itemize}

\begin{tcolorbox}[colback=white, colframe=azul_unam, 
title=Pseudocódigo: cliente.c, breakable]

\textbf{Algoritmo} Cliente

% ===================== DEFINICIONES =====================

\textbf{Constantes y Variables Globales}

\begin{algorithmic}[1]

\State \textbf{Definir} $IP\_SERVIDOR \gets$ "127.0.0.1"
\State \textbf{Definir} $PUERTO \gets 3000$
\State \textbf{Definir} $TAM\_MAX \gets 1024$

\vspace{0.3cm}

% ===================== MAIN =====================

\Procedure{main}{}

\State $sock \gets 0$
\State Definir $serv\_addr$
\State Definir $buffer[TAM\_MAX]$

\State Escribir("[1] Creando socket TCP...")

\If{$(sock \gets socket(AF\_INET, SOCK\_STREAM, 0)) < 0$}
    \State \Return -1
\EndIf

\State $serv\_addr.sin\_family \gets AF\_INET$
\State $serv\_addr.sin\_port \gets htons(PUERTO)$

\If{$inet\_pton(AF\_INET, IP\_SERVIDOR, serv\_addr.sin\_addr) \leq 0$}
    \State \Return -1
\EndIf

\State Escribir("[2] Conectando al servidor...")

\If{$connect(sock, serv\_addr) < 0$}
    \State \Return -1
\EndIf

\State Limpiar(buffer)
\State Recibir$(sock, buffer)$

\State Escribir("\textbackslash n[Mensaje Inicial del Servidor]: ", buffer)

\vspace{0.2cm}

% ===================== BUCLE =====================

\While{Verdadero}

    \State Escribir("========================================")
    \State Escribir(" COMANDOS ACTIVOS:")
    \State Escribir(" CONSULTA | INSERCION | ELIMINACION |")
    \State Escribir(" ACTUALIZACION | OFF ")
    \State Escribir("========================================")
    \State Escribir("Tu mensaje (Comando): ")

    \State Limpiar(buffer)
    \State LeerEntrada(buffer)
    \State EliminarSalto(buffer)

    \If{Buscar(buffer, "OFF") $\neq NULL$}
        \State Escribir("[!] Cerrando cliente...")
        \State Enviar$(sock, "OFF")$
        \State Romper

    \ElsIf{Buscar(buffer, "CONSULTA") $\neq NULL$}
        \State realizar\_consulta$(sock)$

    \ElsIf{Buscar(buffer, "INSERCION") $\neq NULL$}
        \State realizar\_insercion$(sock)$

    \ElsIf{Buscar(buffer, "ELIMINACION") $\neq NULL$}
        \State realizar\_eliminacion$(sock)$

    \ElsIf{Buscar(buffer, "ACTUALIZACION") $\neq NULL$}
        \State realizar\_actualizacion$(sock)$

    \Else
        \State Escribir("[-] Comando no reconocido. Escribe: CONSULTA, INSERCION, ACTUALIZACION, ELIMINACION o OFF.")
    \EndIf

\EndWhile

\State Cerrar$(sock)$
\State Escribir("\textbackslash nDesconectado exitosamente.")

\State \Return 0

\EndProcedure

\vspace{0.3cm}

% ===================== CONSULTA =====================

\Procedure{realizar\_consulta}{sock}

\State Definir $tabla, llave, mascara, paquete\_envio, respuesta\_servidor$

\State Escribir("\textbackslash n\textbackslash t--- MODO CONSULTA ---")
\State Escribir("[ESTUDIANTE,0]\textbackslash t\textbackslash t[DIRECCION,1]")
\State Escribir("[CARRERA,2]\textbackslash t\textbackslash t[HISTORIAL,3]")
\State Escribir("[INSCRIPCION,4]\textbackslash t\textbackslash t[SECCION,5]")
\State Escribir("[PROFESOR,6]\textbackslash t\textbackslash t[DEPARTAMENTO,7]")
\State Escribir("[NIVELES,8]\textbackslash t\textbackslash t[HORARIO,9]")
\State Escribir("[GRADO,10]\textbackslash t\textbackslash t[CURSO,11]")
\State Escribir("[AÑOS,12]\textbackslash t\textbackslash t[SEMESTRE,13]")

\State Escribir("Ingresa el numero de tabla (0-13): ")
\State LeerEntrada(tabla)

\State Escribir("Ingresa la llave a buscar ('full', 'low' o ID): ")
\State LeerEntrada(llave)

\State Escribir("Ingresa la mascara de parametros (Bits 1/0): ")
\State LeerEntrada(mascara)

\State $paquete\_envio \gets "1|" + tabla + "|" + llave + "|" + mascara$
\State Enviar$(sock, paquete\_envio)$

\State Limpiar(respuesta\_servidor)
\State Recibir$(sock, respuesta\_servidor)$

\If{EmpiezaCon(respuesta\_servidor, "CANTIDAD|")}

    \State $total\_filas \gets$ ConvertirEntero$(respuesta\_servidor + 9)$

    \If{Buscar(llave,"full") $\neq NULL$ \textbf{o} Buscar(llave,"low") $\neq NULL$}
        \State $total\_filas \gets total\_filas - 1$
    \EndIf

    \State Escribir("\textbackslash n--- RESULTADOS DE LA BD (", total\_filas, " registros) ---")

    \State $archivo \gets$ Abrir("./consulta.csv","w")

    \State Enviar$(sock,"ACK")$

    \For{$i \gets 0$ to $total\_filas - 1$}

        \State Limpiar(respuesta\_servidor)
        \State Recibir$(sock, respuesta\_servidor)$

        \State Escribir("[", i, "]: ", respuesta\_servidor)
        \State EscribirArchivo$(archivo, respuesta\_servidor)$

        \State Enviar$(sock,"ACK")$

    \EndFor

    \State Escribir("----------------------------------------")
    \State Cerrar$(archivo)$

\Else

    \State Escribir("\textbackslash n[Servidor]: ", respuesta\_servidor)

\EndIf

\EndProcedure

\vspace{0.3cm}

% ===================== INSERCION =====================

\Procedure{realizar\_insercion}{sock}

\State Definir $tabla, paquete\_envio, respuesta\_servidor$
\State Definir $datos\_csv \gets ""$
\State Definir $entrada\_temporal$

\State Escribir("\textbackslash n--- MODO INSERCION ---")
\State Escribir("Ingresa el numero de tabla (0-13): ")
\State LeerEntrada(tabla)

\State $paquete\_envio \gets "2|" + tabla$
\State Enviar$(sock, paquete\_envio)$

\State Limpiar(respuesta\_servidor)
\State Recibir$(sock, respuesta\_servidor)$

\If{EmpiezaCon(respuesta\_servidor,"COLUMNAS|")}

    \State $num\_columnas \gets$ ConvertirEntero$(respuesta\_servidor + 9)$

    \State Escribir("[*] El servidor verifico la tabla. Se requieren ", num\_columnas, " columnas.")
    \State Escribir("--------------------------------------------------")

    \State $datos\_csv \gets "DATOS|"$

    \For{$i \gets 0$ to $num\_columnas - 1$}

        \If{$i = 0$}
            \State Escribir("Ingrese el dato para la columna ", i+1, " (LLAVE PRIMARIA): ")
        \Else
            \State Escribir("Ingrese el dato para la columna ", i+1, ": ")
        \EndIf

        \State LeerEntrada(entrada\_temporal)

        \State Concatenar$(datos\_csv, entrada\_temporal)$

        \If{$i < num\_columnas - 1$}
            \State Concatenar$(datos\_csv, ",")$
        \EndIf

    \EndFor

    \State Escribir("\textbackslash n[Enviando paquete]: ", datos\_csv)
    \State Enviar$(sock, datos\_csv)$

    \State Limpiar(respuesta\_servidor)
    \State Recibir$(sock, respuesta\_servidor)$

    \State Escribir("[Respuesta del Servidor]: ", respuesta\_servidor)

\Else

    \State Escribir("\textbackslash n[Servidor rechazo la peticion]: ", respuesta\_servidor)

\EndIf

\EndProcedure

\vspace{0.3cm}

% ===================== ELIMINACION =====================

\Procedure{realizar\_eliminacion}{sock}

\State Definir $tabla, llave, paquete\_envio, respuesta\_servidor$

\State Escribir("\textbackslash n--- MODO ELIMINACION ---")
\State Escribir("Ingresa el numero de tabla (0-13): ")
\State LeerEntrada(tabla)

\State Escribir("Ingresa el ID (llave primaria) a eliminar: ")
\State LeerEntrada(llave)

\State $paquete\_envio \gets "3|" + tabla + "|" + llave$
\State Enviar$(sock, paquete\_envio)$

\State Limpiar(respuesta\_servidor)
\State Recibir$(sock, respuesta\_servidor)$

\State Escribir("\textbackslash n[Respuesta del Servidor]: ", respuesta\_servidor)

\EndProcedure

\vspace{0.3cm}

% ===================== ACTUALIZACION =====================

\Procedure{realizar\_actualizacion}{sock}

\State Definir $tabla, llave\_vieja, paquete\_envio, respuesta\_servidor$
\State Definir $datos\_csv \gets ""$
\State Definir $entrada\_temporal$

\State Escribir("\textbackslash n--- MODO ACTUALIZACION ---")
\State Escribir("Ingresa el numero de tabla (0-13): ")
\State LeerEntrada(tabla)

\State Escribir("Ingresa la llave ORIGINAL del registro a modificar: ")
\State LeerEntrada(llave\_vieja)

\State $paquete\_envio \gets "4|" + tabla + "|" + llave\_vieja$
\State Enviar$(sock, paquete\_envio)$

\State Limpiar(respuesta\_servidor)
\State Recibir$(sock, respuesta\_servidor)$

\If{EmpiezaCon(respuesta\_servidor,"COLUMNAS|")}

    \State $num\_columnas \gets$ ConvertirEntero$(respuesta\_servidor + 9)$

    \State Escribir("[*] Registro encontrado. Ingresa los nuevos datos (", num\_columnas, " columnas).")
    \State Escribir("--------------------------------------------------")

    \State $datos\_csv \gets "DATOS|"$

    \For{$i \gets 0$ to $num\_columnas - 1$}

        \If{$i = 0$}
            \State Escribir("NUEVA Llave Primaria (Puede ser la misma '", llave\_vieja, "'): ")
        \Else
            \State Escribir("NUEVO dato para la columna ", i+1, ": ")
        \EndIf

        \State LeerEntrada(entrada\_temporal)

        \State Concatenar$(datos\_csv, entrada\_temporal)$

        \If{$i < num\_columnas - 1$}
            \State Concatenar$(datos\_csv, ",")$
        \EndIf

    \EndFor

    \State Escribir("\textbackslash n[Enviando actualizacion]: ", datos\_csv)
    \State Enviar$(sock, datos\_csv)$

    \State Limpiar(respuesta\_servidor)
    \State Recibir$(sock, respuesta\_servidor)$

    \State Escribir("[Respuesta del Servidor]: ", respuesta\_servidor)

\Else

    \State Escribir("\textbackslash n[Servidor rechazo la peticion]: ", respuesta\_servidor)

\EndIf

\EndProcedure

\end{algorithmic}
\end{tcolorbox}

%----- PSEUDOCODIGO SERVIDOR.C ------

\begin{tcolorbox}[colback=white, colframe=azul_unam, 
title=Pseudocódigo: servidor.c, breakable]

\textbf{Algoritmo} Servidor

% ===================== DEFINICIONES =====================

\textbf{Constantes y Variables Globales}

\begin{algorithmic}[1]

\State \textbf{Definir} $PUERTO \gets 3000$

\State \textbf{Definir} $directorio \in \text{ARCHIVERO}$

\State \textbf{Definir} $escritura \gets$ ``./buffer/''
\State \textbf{Definir} $carpeta \gets$ ``./Datos/''
\State \textbf{Definir} $extension \gets$ ``.csv''

\State \textbf{Definir} $archivos \gets$ \{
``Estudiante'', ``Direccion'', ``Carrera'', ``Historial'',
``Inscripcion'', ``Seccion'', ``Profesor'', ``Departamento'',
``Niveles'', ``Horario'', ``Grado'', ``Curso'', ``Años'', ``Semestre''
\}

\State \textbf{Definir} $tamano \gets$ Tamaño$(archivos)$

\State \textbf{Definir} $TAM\_MAX \gets 1024$

\vspace{0.3cm}

% ===================== MAIN =====================

\Procedure{main}{}

\State Crear hilo\_directorio
\State Crear hilo\_red

\State datos\_dir $\gets$ \{directorio, carpeta, escritura, archivos, extension, tamano\}
\State datos\_red $\gets$ \{PUERTO, -1\}

\State Escribir("[Servidor] Iniciando hilos de preparacion...")

\State CrearHilo$(hilo\_directorio, funcion\_h1, datos\_dir)$
\State CrearHilo$(hilo\_red, funcion\_h2, datos\_red)$

\State EsperarHilo$(hilo\_directorio)$
\State EsperarHilo$(hilo\_red)$

\State $server\_fd \gets datos\_red.socket\_fd\_resultado$

\State Escribir("[+] Servidor iniciado. Esperando conexion...")

\State $nuevo\_socket \gets$ AceptarConexion$(server\_fd)$

\If{$nuevo\_socket < 0$}
    \State Error("ERROR: Falla al intentar conectar")
    \State Cerrar$(server\_fd)$
    \State \Return
\EndIf

\State Escribir("[+] Cliente conectado. Iniciando sesion única.")
\State $Enviar$(nuevo\_socket, "Conexion establecida, el servidor finalizara despues de la sesion.\textbackslash n")

\vspace{0.2cm}

% ===================== LOOP =====================

\While{Verdadero}

    \State Limpiar$(buffer)$
    \State $bytes\_leidos \gets$ Recibir$(nuevo\_socket, buffer)$

    \If{$bytes\_leidos \leq 0 \lor \text{CompararIgnoreCase(buffer, "OFF")} = 0$}
        \State Escribir("Cliente desconectado o se recibio comando OFF. Finalizando...")
        \State Romper
    \EndIf

    \State $op\_str \gets$ Token$(buffer, "|")$
    \If{$op\_str = NULL$}
        \State Continuar
    \EndIf

    \State $operacion \gets$ ConvertirEntero$(op\_str)$

    \State $tabla\_str \gets$ Token$(NULL, "|")$
    \If{$tabla\_str = NULL$}
        \State Continuar
    \EndIf

    \State $tabla \gets$ ConvertirEntero$(tabla\_str)$

    % ===================== CONSULTA =====================

    \If{$operacion = 1$}
        \State $llave \gets$ Token$(NULL, "|")$
        \State $mascara \gets$ Token$(NULL, RESTO)$
        \State procesar\_consulta$(nuevo\_socket, tabla, llave, mascara, directorio)$

    % ===================== INSERCION =====================

    \ElsIf{$operacion = 2$}

        \State $ruta \gets directorio.rutas[tabla]$
        \State $analisis \gets analizar\_archivo(ruta)$

        \State Enviar(nuevo\_socket, "COLUMNAS\textbar" + ConvertirCadena(analisis.num\_columnas))
        \State Limpiar(buffer)
        \State bytes\_datos $\gets$ Recibir(nuevo\_socket, buffer)

        \If{bytes\_datos > 0 \textbf{Y} EmpiezaCon(buffer, "DATOS|")}

            \State $datos\_csv \gets buffer + 6$

            \State Crear peticion\_ins
            \State peticion\_ins.numero\_tabla $\gets tabla$
            \State peticion\_ins.datos $\gets datos\_csv$
            \State peticion\_ins.error $\gets NULL$

            \State solicitud\_insercion$(peticion\_ins, directorio)$

            \If{$peticion\_ins.error \neq NULL$}
                \State Enviar$(nuevo\_socket, peticion\_ins.error)$
                \State Liberar$(peticion\_ins.error)$
            \Else
                \State $Enviar$(nuevo\_socket, "ERROR: Fallo desconocido en insercion.\textbackslash n")
            \EndIf

        \Else
            \State $Enviar$(nuevo\_socket, "ERROR: Se rompio el protocolo de insercion.\textbackslash n")
        \EndIf

    % ===================== ELIMINACION =====================

    \ElsIf{$operacion = 3$}

        \State $llave \gets$ Token$(NULL, RESTO)$

        \If{$llave \neq NULL$}

            \State Crear peticion\_del
            \State peticion\_del.num\_tabla $\gets tabla$
            \State peticion\_del.primary\_key $\gets llave$

            \State solicitud\_eliminacion$(peticion\_del, directorio)$

            \If{$peticion\_del.error \neq NULL$}
                \State Enviar$(nuevo\_socket, peticion\_del.error)$
                \State Liberar$(peticion\_del.error)$
            \Else
                \State $Enviar$(nuevo\_socket, "ERROR: Fallo desconocido en eliminacion.\textbackslash n")
            \EndIf

        \EndIf

    % ===================== UPDATE =====================

    \ElsIf{$operacion = 4$}

        \State $llave\_vieja \gets$ Token$(NULL, RESTO)$

        \State $ruta \gets directorio.rutas[tabla]$
        \State $estado \gets validar\_llave(llave\_vieja, tabla, ruta)$

        \If{Comparar$(estado, "EXITO") = 0$}
            \State $Enviar$(nuevo\_socket, "ERROR: El registro que intentas actualizar no existe.\textbackslash n")
            \State Continuar
        \EndIf

        \State Copiar$(llave\_respaldo, llave\_vieja)$

        \State analisis $\gets$ analizar\_archivo(ruta)
        \State Enviar(nuevo\_socket, "COLUMNAS\textbar" + ConvertirCadena(analisis.num\_columnas))

        \State Limpiar(buffer)
        \State bytes\_datos $\gets$ Recibir(nuevo\_socket, buffer)

        \If{bytes\_datos > 0 \textbf{Y} EmpiezaCon(buffer, "DATOS|")}

            \State $datos\_csv \gets buffer + 6$

            \State Crear peticion\_upd
            \State peticion\_upd.num\_tabla $\gets tabla$
            \State peticion\_upd.primary\_key $\gets llave\_respaldo$
            \State peticion\_upd.parametros $\gets datos\_csv$

            \State solicitud\_update$(peticion\_upd, directorio)$

            \If{$peticion\_upd.error \neq NULL$}
                \State Enviar$(nuevo\_socket, peticion\_upd.error)$
                \State Liberar$(peticion\_upd.error)$
            \Else
                \State $Enviar$(nuevo\_socket, "EXITO: El registro fue actualizado correctamente en el servidor.\textbackslash n")
            \EndIf

        \Else
            \State $Enviar$(nuevo\_socket, "ERROR: Se rompio el protocolo de actualizacion.\textbackslash n")
        \EndIf

    \EndIf

\EndWhile

\State Cerrar$(nuevo\_socket)$
\State Cerrar$(server\_fd)$

\State Escribir("[*] Servidor apagado...")

\EndProcedure

\vspace{0.3cm}

% ===================== HILOS =====================

\Procedure{funcion\_h1}{args}
\State crear\_directorio$(args)$
\EndProcedure

\Procedure{funcion\_h2}{args}
\State args.socket\_fd\_resultado $\gets$ levantar\_servicio$(args.puerto)$
\EndProcedure

\end{algorithmic}
\end{tcolorbox}

%--------- PSEUDOCODIGO FUNCIONES.C ----------

\begin{tcolorbox}[colback=white, colframe=azul_unam, 
title=Pseudocódigo: funciones.c, breakable]

\textbf{Algoritmo} Funciones

\begin{algorithmic}[1]

% ===================== CREAR DIRECTORIO =====================

\Procedure{crear\_directorio}{dir, carpeta, archivos, tipo, update, cantidad\_archivos}

\State Escribir("Se esta creando un directorio con ", cantidad\_archivos, " rutas de archivos ", tipo)

\For{$i \gets 0$ to $cantidad\_archivos - 1$}

    \State $len\_carpeta \gets$ Longitud$(carpeta)$ + Longitud$(archivos[i])$ + Longitud$(tipo)$ + 1
    \State $len\_update \gets$ Longitud$(update)$ + Longitud$(archivos[i])$ + Longitud$(tipo)$ + 1

    \State $ruta\_completa \gets$ malloc$(len\_carpeta)$
    \State $ruta\_updates \gets$ malloc$(len\_update)$

    \If{$ruta\_completa = NULL \lor ruta\_updates = NULL$}
        \State Escribir("ERROR: Memoria insuficiente para crear rutas.")
        \State \Return
    \EndIf

    \State Copiar$(ruta\_completa, carpeta)$
    \State Concatenar$(ruta\_completa, archivos[i])$
    \State Concatenar$(ruta\_completa, tipo)$

    \State Copiar$(ruta\_updates, update)$
    \State Concatenar$(ruta\_updates, archivos[i])$
    \State Concatenar$(ruta\_updates, tipo)$

    \State $dir.rutas[i] \gets ruta\_completa$
    \State $dir.rutas\_update[i] \gets ruta\_updates$
    \State $dir.cantidad \gets i + 1$

\EndFor

\EndProcedure

\vspace{0.3cm}

% ===================== ANALIZAR ARCHIVO =====================

\Procedure{analizar\_archivo}{ruta}

\State $analisis \gets \{0,0\}$
\State $archivo \gets$ Abrir$(ruta)$

\State $primera\_linea \gets 1$

\While{LeerLinea$(archivo, linea)$}

    \State $analisis.num\_lineas \gets analisis.num\_lineas + 1$

    \If{$primera\_linea = 1$}

        \State $ptr \gets linea$

        \While{$ptr \neq FIN$}
            \If{$ptr = ','$}
                \State $analisis.num\_columnas \gets analisis.num\_columnas + 1$
            \EndIf
            \State Avanzar$(ptr)$
        \EndWhile

        \State $analisis.num\_columnas \gets analisis.num\_columnas + 1$
        \State $primera\_linea \gets 0$

    \EndIf

\EndWhile

\State Cerrar$(archivo)$
\State \Return $analisis$

\EndProcedure

\vspace{0.3cm}

% ===================== VALIDAR LLAVE =====================

\Procedure{validar\_llave}{llave, num\_table, ruta}

\State $analisis \gets analizar\_archivo(ruta)$
\State $columnas \gets analisis.num\_columnas$

\State $mascara \gets$ malloc$(columnas + 1)$

\If{$mascara = NULL$}
    \State \Return "ERROR: Memoria insuficiente"
\EndIf

\State $mascara[0] \gets '1'$

\For{$i \gets 1$ to $columnas - 1$}
    \State $mascara[i] \gets '0'$
\EndFor

\State $mascara[columnas] \gets$ '\textbackslash 0'

\State $tmp.numero\_tabla \gets num\_table$
\State $tmp.error \gets NULL$
\State $tmp.llave \gets llave$
\State $tmp.cantidad\_resultados \gets 0$
\State $tmp.resultado \gets NULL$
\State $tmp.parametros \gets mascara$

\State consulta\_tabla$(tmp, directorio, columnas)$

\State Liberar$(mascara)$

\If{$tmp.cantidad\_resultados > 0$}

    \For{$i \gets 0$ to $tmp.cantidad\_resultados - 1$}
        \State Liberar$(tmp.resultado[i])$
    \EndFor

    \State Liberar$(tmp.resultado)$

    \State \Return "ERROR: La llave se repite en la tabla\textbackslash n"

\EndIf

\State \Return "EXITO"

\EndProcedure

\vspace{0.3cm}

% ===================== LEVANTAR SERVICIO =====================

\Procedure{levantar\_servicio}{puerto}

\State $servidor\_fd \gets socket(AF\_INET, SOCK\_STREAM)$

\If{$servidor\_fd = 0$}
    \State Error("ERROR: Fallo al crear el socket TCP")
    \State Terminar
\EndIf

\State setsockopt$(servidor\_fd, SO\_REUSEADDR)$

\State Limpiar$(direct)$
\State $direct.sin\_family \gets AF\_INET$
\State $direct.sin\_addr \gets INADDR\_ANY$
\State $direct.sin\_port \gets htons(puerto)$

\If{bind$(servidor\_fd, direct) < 0$}
    \State Error("ERROR: Fallo en bind (¿Puerto en uso?)")
    \State Terminar
\EndIf

\If{listen$(servidor\_fd, 5) < 0$}
    \State Error("ERROR: Fallo en listen")
    \State Terminar
\EndIf

\State Escribir("[+] Servicio TCP levantado. Escuchando en el puerto ", puerto)

\State \Return $servidor\_fd$

\EndProcedure

\end{algorithmic}
\end{tcolorbox}

% ------- PSEUDOCODIGO CONSULTA -------

\begin{tcolorbox}[colback=white, colframe=azul_unam, 
title=Pseudocódigo: consulta.c, breakable]

\textbf{Algoritmo} Consulta

\begin{algorithmic}[1]

% ===================== PROCESAR CONSULTA =====================

\Procedure{procesar\_consulta}{socket\_cliente, tabla, llave, mascara, dir}

\If{$llave = NULL \lor mascara = NULL$}
    \State $msj \gets$ ``ERROR: Faltan parametros en la consulta.\textbackslash n''
    \State Enviar$(socket\_cliente, msj)$
    \State \Return
\EndIf

\State $c \gets \{tabla, NULL, mascara, llave, 0, NULL\}$
\State solicitud\_consulta$(c, dir)$

\If{$c.error \neq NULL$}
    \If{Subcadena$(c.error,0,5) = $ ``ERROR''}
        \State Enviar$(socket\_cliente, c.error)$
        \State Liberar$(c.error)$
        \State \Return
    \EndIf
    \State Liberar$(c.error)$
\EndIf

\If{$c.cantidad\_resultados = 0$}
    \State $msj \gets$ ``No se encontraron registros coincidentes.\textbackslash n''
    \State Enviar$(socket\_cliente, msj)$
    \State \Return
\EndIf

\State $metadatos \gets$ ``CANTIDAD|'' $+$ Convertir\_Cadena$(c.cantidad\_resultados)$
\State Enviar$(socket\_cliente, metadatos)$

\State Recibir$(socket\_cliente, buffer\_ack)$

\For{$i \gets 0$ to $c.cantidad\_resultados - 1$}
    \State Enviar$(socket\_cliente, c.resultado[i])$
    \State Recibir$(socket\_cliente, buffer\_ack)$
    \State Liberar$(c.resultado[i])$
\EndFor

\State Liberar$(c.resultado)$

\EndProcedure

\vspace{0.3cm}

% ===================== SOLICITUD CONSULTA =====================

\Procedure{solicitud\_consulta}{cliente, dir}

\If{$cliente.numero\_tabla < 0 \lor cliente.numero\_tabla \geq dir.cantidad$}
    \State $cliente.error \gets$ ``ERROR: LA TABLA SOLICITADA NO EXISTE\textbackslash n''
    \State \Return
\EndIf

\State $ruta \gets dir.rutas[cliente.numero\_tabla]$
\State $analisis \gets analizar\_archivo(ruta)$
\State $columnas\_reales \gets analisis.num\_columnas$

\If{Longitud$(cliente.parametros) \neq columnas\_reales$}
    \State $cliente.error \gets$ ``ERROR: La mascara ingresada no coincide. Esta tabla requiere exactamente'' 
    \State $cliente.error \gets$ cliente.error + Convertir\_Cadena(columnas\_reales) + ``parametros (Bits 1/0).''
\Else
    \State consulta\_tabla$(cliente, dir, columnas\_reales)$
\EndIf

\EndProcedure

\vspace{0.3cm}

% ===================== CONSULTA TABLA =====================

\Procedure{consulta\_tabla}{cliente, dir, num\_columnas}

\State $analisis \gets analizar\_archivo(dir.rutas[cliente.numero\_tabla])$
\State $archivo \gets$ Abrir$(dir.rutas[cliente.numero\_tabla])$

\State $hallados \gets 0$
\State $cliente.resultado \gets$ Arreglo$(analisis.num\_lineas)$

\While{LeerLinea$(archivo, linea)$}

    \State $linea \gets$ EliminarSalto$(linea)$

    \State $tokens \gets$ Separar$(linea, ",")$
    \State $col \gets 0$

    \For{$i \gets 0$ \textbf{mientras} $tokens[i] \neq NULL \land col < num\_columnas$}
        \State $col \gets col + 1$
    \EndFor

    \If{$col = 0$}
        \State \textbf{continuar}
    \EndIf

    \State $es\_full \gets (cliente.llave \neq NULL \land cliente.llave = ``full'')$
    \State $es\_low \gets (cliente.llave \neq NULL \land cliente.llave = ``low'')$
    \State $match \gets 0$

    \If{$es\_full \lor es\_low$}
        \State $match \gets 1$
    \ElsIf{$cliente.llave = NULL$}
        \State $match \gets 1$
    \ElsIf{$tokens[0] \neq NULL \land tokens[0] = cliente.llave$}
        \State $match \gets 1$
    \EndIf

    \If{$match = 0$}
        \State \textbf{continuar}
    \EndIf

    \State $buffer \gets ""$
    \State $primero \gets 1$

    \If{$es\_full$}
        \For{$i \gets 0$ to $col - 1$}
            \If{$primero = 0$}
                \State $buffer \gets buffer + ","$
            \EndIf
            \State $buffer \gets buffer + tokens[i]$
            \State $primero \gets 0$
        \EndFor
    \Else
        \For{$i \gets 0$ to $num\_columnas - 1$}
            \If{$i < col$}
                \If{$cliente.parametros[i] = '1'$}
                    \If{$primero = 0$}
                        \State $buffer \gets buffer + ","$
                    \EndIf
                    \State $buffer \gets buffer + tokens[i]$
                    \State $primero \gets 0$
                \EndIf
            \EndIf
        \EndFor
    \EndIf

    \State $cliente.resultado[hallados] \gets buffer$
    \State $hallados \gets hallados + 1$

\EndWhile

\State $cliente.cantidad\_resultados \gets hallados$

\If{$hallados > 0 \land hallados < analisis.num\_lineas$}
    \State Redimensionar$(cliente.resultado, hallados)$
\EndIf

\If{$hallados = 0$}
    \State Liberar$(cliente.resultado)$
    \State $cliente.resultado \gets NULL$
    \State $cliente.error \gets$ ``ERROR: No se encontro ningun registro en la tabla consultada\textbackslash n''
\Else
    \State $cliente.error \gets$ ``EXITO: Se encontro al menos 1 registro en la tabla\textbackslash n''
\EndIf

\State Cerrar$(archivo)$

\EndProcedure

\end{algorithmic}
\end{tcolorbox}

% ----------- PSEUDOCODIGO INSERCION.H ------------
\begin{tcolorbox}[colback=white, colframe=azul_unam, 
title=Pseudocódigo: insercion.c, breakable]

\textbf{Algoritmo} Insercion y Eliminacion

\begin{algorithmic}[1]

% ===================== SOLICITUD INSERCION =====================

\Procedure{solicitud\_insercion}{cliente, dir}

\If{$cliente.numero\_tabla < 0 \lor cliente.numero\_tabla > 13$}
    \State $cliente.error \gets$ ``ERROR: Tabla no valida para insercion''
    \State \Return
\EndIf

\State $cond \gets$ Buscar$(cliente.datos, ",,") \neq NULL$
\State $cond \gets cond \lor \text{cliente.datos[0] = ','}$
\State $cond \gets cond \lor \text{cliente.datos[Longitud(cliente.datos)-1] = ','}$

\If{$cond$}
    \State $cliente.error \gets$ ``ERROR: Los datos contienen campos vacios.''
    \State \Return
\EndIf

\State Copiar$(datos\_copia, cliente.datos)$
\State $llave\_nueva \gets$ Token$(datos\_copia, ",")$

\If{$llave\_nueva = NULL$}
    \State $cliente.error \gets$ ``ERROR: Formato de datos invalido.''
    \State \Return
\EndIf

\State $ruta\_archivo \gets dir.rutas[cliente.numero\_tabla]$

\State $estado \gets$ validar\_llave$(llave\_nueva, cliente.numero\_tabla, ruta\_archivo)$

\If{Comparar$(estado, "EXITO") \neq 0$}
    \State $cliente.error \gets$ ``ERROR: La llave primaria ya existe en la base de datos.''
    \State \Return
\EndIf

\State $archivo \gets$ Abrir$(ruta\_archivo, "a")$

\If{$archivo = NULL$}
    \State $cliente.error \gets$ ``ERROR: No se pudo abrir el archivo para insertar.''
    \State \Return
\EndIf

\State $EscribirArchivo$(archivo, cliente.datos + "\textbackslash n")
\State Cerrar$(archivo)$

\State $cliente.error \gets$ ``EXITO: Registro insertado correctamente.''

\EndProcedure

\vspace{0.3cm}

% ===================== SOLICITUD ELIMINACION =====================

\Procedure{solicitud\_eliminacion}{cliente, dir}

\State argumentos $\gets$ cliente.primary\_key, cliente.num\_tabla, dir.rutas[cliente.num\_tabla]
\State estado $\gets$ validar\_llaves(argumentos)

\If{Comparar$(estado, "EXITO") = 0$}
    \State $cliente.error \gets$ ``ERROR: La llave no existe. No se puede eliminar un registro fantasma.''
    \State \Return
\EndIf

\State $ruta\_original \gets dir.rutas[cliente.num\_tabla]$
\State $ruta\_temporal \gets dir.rutas\_update[cliente.num\_tabla]$

\State $archivo\_lectura \gets$ Abrir$(ruta\_original, "r")$
\State $archivo\_escritura \gets$ Abrir$(ruta\_temporal, "w")$

\If{$archivo\_lectura = NULL \lor archivo\_escritura = NULL$}
    \State $cliente.error \gets$ ``ERROR: Fallo critico al abrir los archivos para eliminacion.''
    \If{$archivo\_lectura \neq NULL$}
        \State Cerrar$(archivo\_lectura)$
    \EndIf
    \If{$archivo\_escritura \neq NULL$}
        \State Cerrar$(archivo\_escritura)$
    \EndIf
    \State \Return
\EndIf

\While{LeerLinea$(archivo\_lectura, linea)$}

    \State Copiar$(linea\_copia, linea)$
    \State $llave\_actual \gets$ Token$(linea\_copia, ",")$

    \If{$llave\_actual \neq NULL \land Comparar$(llave\_actual, cliente.primary\_key$) = 0$}
        \State \textbf{continuar}
    \Else
        \State EscribirArchivo$(archivo\_escritura, linea)$
    \EndIf

\EndWhile

\State Cerrar$(archivo\_lectura)$
\State Cerrar$(archivo\_escritura)$

\State EliminarArchivo$(ruta\_original)$

\If{Renombrar$(ruta\_temporal, ruta\_original) = 0$}
    \State $cliente.error \gets$ ``EXITO: Registro eliminado permanentemente.''
\Else
    \State $cliente.error \gets$ ``ERROR: No se pudo completar la eliminacion fisica.''
\EndIf

\EndProcedure

\end{algorithmic}
\end{tcolorbox}

% ---------- PSEUDOCODIGO UPDATE.C -----------

\begin{tcolorbox}[colback=white, colframe=azul_unam, 
title=Pseudocódigo: update.c, breakable]

\textbf{Algoritmo} Update

\begin{algorithmic}[1]

% ===================== SOLICITUD UPDATE =====================

\Procedure{solicitud\_update}{cliente, dir}

\State argumentos $\gets$ cliente.primary\_key, cliente.num\_tabla, dir.rutas[cliente.num\_tabla]
\State estado\_origen $\gets$ validar\_llaves(argumentos)

\If{Comparar$(estado\_origen, "EXITO") = 0$}
    \State $cliente.error \gets$ ``ERROR: El registro que intentas actualizar no existe.\textbackslash n''
    \State \Return
\EndIf

\State Copiar$(parametros\_copia, cliente.parametros)$
\State $nueva\_llave \gets$ Token$(parametros\_copia, ",")$

\If{$nueva\_llave \neq NULL \land Comparar$(cliente.primary\_key, nueva\_llave$) \neq 0$}

    \State args\_destino $\gets$ nueva\_llave, cliente.num\_tabla, dir.rutas[cliente.num\_tabla]
    \State estado\_destino $\gets$ validar\_llave(args\_destino)

    \If{Comparar$(estado\_destino, "EXITO") \neq 0$}
        \State $cliente.error \gets$ ``ERROR: La nueva llave primaria ya esta en uso por otro registro.\textbackslash n''
        \State \Return
    \EndIf

\EndIf

\State $cliente.error \gets NULL$
\State update\_archivo$(cliente, dir)$

\EndProcedure

\vspace{0.3cm}

% ===================== UPDATE ARCHIVO =====================

\Procedure{update\_archivo}{cliente, dir}

\State $ruta\_original \gets dir.rutas[cliente.num\_tabla]$
\State $ruta\_temporal \gets dir.rutas\_update[cliente.num\_tabla]$

\State $archivo\_lectura \gets$ Abrir$(ruta\_original, "r")$
\State $archivo\_escritura \gets$ Abrir$(ruta\_temporal, "w")$

\If{$archivo\_lectura = NULL \lor archivo\_escritura = NULL$}
    \State $cliente.error \gets$ ``ERROR: Fallo critico al abrir/crear los archivos.\textbackslash n''
    \If{$archivo\_lectura \neq NULL$}
        \State Cerrar$(archivo\_lectura)$
    \EndIf
    \If{$archivo\_escritura \neq NULL$}
        \State Cerrar$(archivo\_escritura)$
    \EndIf
    \State \Return
\EndIf

\While{LeerLinea$(archivo\_lectura, linea)$}

    \State Copiar$(linea\_copia, linea)$
    \State $llave\_actual \gets$ Token$(linea\_copia, ",")$

    \If{$llave\_actual \neq NULL \land Comparar$(llave\_actual, cliente.primary\_key$) = 0$}
        \State $EscribirArchivo$(archivo\_escritura, cliente.parametros + "\textbackslash n")
    \Else
        \State EscribirArchivo$(archivo\_escritura, linea)$
    \EndIf

\EndWhile

\State Cerrar$(archivo\_lectura)$
\State Cerrar$(archivo\_escritura)$

\State EliminarArchivo$(ruta\_original)$

\If{Renombrar$(ruta\_temporal, ruta\_original) \neq 0$}
    \State $cliente.error \gets$ ``ERROR: No se pudo mover el archivo temporal.\textbackslash n''
    \State \Return
\EndIf

\State Escribir("EXITO: Se actualizo el registro ", cliente.primary\_key)

\EndProcedure

\end{algorithmic}
\end{tcolorbox}